\section{The exact zero-risk resource spectrum}\label{sec:zero-spectrum}
The mass of one favorable atom gives a sufficient resource threshold, not in general a necessary one. Other zero-loss atoms may be available to the same decoder. We determine the complete finite-width threshold, for the general acquired loss first and then for squared prediction. This statement is about a genuinely acquired probability law; converting it to a physical budget still requires the exact transfer hypotheses.

Use the compact decision space and nonnegative conditional losses $g_r(x)$ of Section~\ref{sec:loss-transfer}. For a dictionary $C$, put $d_C(x)=\min_{r\in C}g_r(x)$ and define
\[
 z_m=\max_{C\subset D,\,1\le |C|\le m}\nu\{d_C=0\}.
\]
\begin{theorem}[Exact finite-width zero-loss criterion]\label{thm:zero-general}
The maximum defining $z_m$ is attained. For every finite $m\ge1$ and finite $b\ge1$,
\begin{equation}\label{eq:zero-general}
 \Q_m^\ell(b)=0\quad\Longleftrightarrow\quad bz_m\ge1.
\end{equation}
When the inequality holds, an optimal procedure accepts only reports in a zero-loss set of an attaining dictionary. When it fails, the attained profile is strictly positive. Neither implication requires positive residual variance, a dimension assumption, a density, or an absence of atoms.
\end{theorem}
\begin{proof}
Represent a dictionary by a tuple in compact $D^m$, allowing repetitions. Conditional losses are continuous in each decision on one common full-measure set. For a convergent sequence of tuples, the nonnegative minima converge pointwise. The indicator of the closed condition $d_C(x)=0$ is upper semicontinuous in the tuple. Reverse Fatou for bounded functions therefore makes its integral upper semicontinuous; compactness gives a maximizing tuple.

If $v=\nu\{d_C=0\}\ge1/b$, the density $h=v^{-1}\one_{\{d_C=0\}}$ has unit mass and cap at most $b$, with objective zero. Conversely, take an attained optimal tuple and density for a zero profile, whose existence was proved in Theorem~\ref{thm:width-transfer}. Nonnegativity implies $h(x)d_C(x)=0$ almost everywhere. Thus its unit mass is supported on $\{d_C=0\}$ and is at most $b\nu\{d_C=0\}\le bz_m$. If $bz_m<1$, no optimizer has zero value; nonnegativity and attainment then make the value strictly positive.
\end{proof}

For squared prediction, form the finite measure on $H$
\[
 \mu_0(A)=\nu\{x:w(x)=0,\ f(x)\in A\}.
\]
List its positive atomic masses in nonincreasing order as $a_1\ge a_2\ge\cdots$, padding with zeros if there are finitely many, and set $a_j=0$ for every $j$ if there are none. The total mass is at most one. Repeated values in this list belong to distinct atom locations.

\begin{theorem}[Atomic zero-risk spectrum and exact resource thresholds]\label{thm:zero-spectrum}
For every $m\ge1$,
\begin{equation}\label{eq:atomic-zero-spectrum}
 z_m=\sum_{j=1}^m a_j,\qquad
 \Q_m(b)=0\quad\Longleftrightarrow\quad
 b\sum_{j=1}^m a_j\ge1.
\end{equation}
The nonatomic part of $\mu_0$ contributes nothing at a fixed finite width, even when $w=0$ everywhere. All finite-width zero-risk acquisition thresholds are therefore determined by the ordered atomic masses of the zero-residual predictive law, not by its matrix rank or its ambient dimension.

Under the exact matched certificates with identity exports, the checkpoint threshold at width $m$ is $N=\kappa/z_m$ when $z_m>0$, and no finite budget yields zero risk when $z_m=0$. For the event-stopped equality, replace $m$ by the total width $M\ge B$. For the state-stopped equality with a sharp mandatory live count, use total labels $B+m$. These conclusions are necessary and sufficient within those declared exact-transfer classes, not for arbitrary experimental interfaces.
\end{theorem}
\begin{proof}
For squared loss, $d_C=w+\dist(f,C)^2$ is zero exactly when $w=0$ and $f\in C$. Hence $\nu\{d_C=0\}=\mu_0(C)$. A finite set receives mass only from atoms. There are at most countably many positive atoms: for each positive integer $n$, at most $n$ atoms can have mass exceeding $1/n$. A nonempty collection of remaining positive masses has a largest member. Indeed if its supremum is $a>0$, only finitely many members can exceed $a/2$, and one of this finite set attains the supremum. Successive removal therefore defines the displayed ordering. Choosing the locations of its first $m$ positive masses maximizes $\mu_0(C)$ among all dictionaries of size at most $m$. Missing locations may be filled by arbitrary repeated points in $H$. This proves the first identity; Theorem~\ref{thm:zero-general} proves the equivalence. Finally substitute $b=N/\kappa$ into the respective exact transfer equalities, keeping their different memory counts.
\end{proof}

\begin{example}[Why one atom does not determine the necessary threshold]\label{ex:two-zero-atoms}
Let the acquired marked law put mass $1/4$ at $(0,0)$ and mass $3/4$ at $(1,0)$, in a carrier containing $[0,1]$. Treating only the first atom would suggest budget cap four. In fact $z_1=3/4$ and $z_2=1$. Thus $\Q_1(b)=0$ exactly for $b\ge4/3$, while $\Q_2(b)=0$ already at $b=1$. The remainder has zero residual variance, which is precisely why the single-atom necessary threshold does not apply. Proposition~\ref{prop:atom} now explicitly states the additional hypothesis needed for its one-atom necessary threshold; the remainder of that proposition is unchanged.
\end{example}

\subsection{An exact variational form, including ties}
The following fixed-dictionary formula records the classical selection mechanism explicitly. It is a population trimming or bounded-density variational principle, not a new invention of a trimming functional \cite{CGM97,GG99}. Its role here is to keep boundary mass and the physical acquisition cap in the same proof.
\begin{proposition}[Selected-mass dual and optimal selection]\label{prop:selection-dual}
For $0\le d\le L$ measurable on a probability space and $b\ge1$,
\begin{equation}\label{eq:selection-dual}
 \min_{0\le h\le b,\ \int h\,d\nu=1}\int hd\,d\nu
 =\max_{0\le t\le L}\left(t-b\int(t-d)_+\,d\nu\right).
\end{equation}
There is a threshold $t$ satisfying $\nu(d<t)\le1/b\le\nu(d\le t)$. An optimizer has density $b$ on $\{d<t\}$, zero on $\{d>t\}$ and a constant in $[0,b]$ on the tie set, chosen to give unit mass. More generally these saturation rules characterize every optimum for an attaining threshold, up to null sets. Applying this formula with $d=d_C$ and then minimizing over $C\in D^m$ gives the exact acquired profile; no interchange of the dictionary minimum and threshold maximum is asserted.
\end{proposition}
\begin{proof}
For every feasible $h$ and every $t$, pointwise minimization of $h(d-t)$ over $[0,b]$ gives $h(d-t)\ge-b(t-d)_+$. Integrating and using unit mass proves weak duality. Take a lower $1/b$ quantile of $d$ in $[0,L]$; elementary continuity of probability measures from below and above gives the two quantile inequalities, including $b=1$ and atoms at the endpoints. The prescribed density exists because the missing mass lies between zero and $b\nu(d=t)$. It attains equality in the pointwise inequality, proving equality and both attainments. The nonnegative difference in that inequality is $h(d-t)$ on $d>t$ and $(b-h)(t-d)$ on $d<t$. Its integral vanishes exactly under the stated saturation rules. On $d=t$ it is zero for every allowed density. This also proves the claim about all optimizers and ties.
\end{proof}
